\section{Amplitude floors}
\label{sec:amplitude}

Theorem~\ref{thm:main} prices a course of treatment. It says nothing about a
course of unbounded length and unbounded cost, because its right side tends to
zero as $B\to\infty$. The second obstruction concerns exactly that regime: an
actuator that is too weak cannot eradicate however long it is used.

\begin{theorem}[Amplitude floor]\label{thm:amp}
Let $\mathcal V\subseteq[0,v_{\max}]$ be the admitted action set and suppose
$q\in(0,1]^m$ satisfies
\begin{equation}\label{eq:ampcert}
 \Phi_v(q)\le0\qquad\text{for every }v\in\mathcal V .
\end{equation}
Then for every predictable policy with values in $\mathcal V$, with no budget
restriction, every initial configuration $z$ and every $t\ge0$,
\begin{equation}\label{eq:ampbound}
 \Prob_z(|Z_t|>0)\ \ge\ 1-\prod_iq_i^{z_i},
\end{equation}
and the same bound holds for the probability of eventual non-extinction, and for
eventual non-extinction after any continuation with
$\Phi_{\mathrm{off}}(q)\le0$. Because $v\mapsto\Phi_v(q)$ is affine,
\eqref{eq:ampcert} for $\mathcal V=[0,v_{\max}]$ is equivalent to the two
inequalities $\PhiB(q)\le0$ and $\Phi_{v_{\max}}(q)\le0$.
\end{theorem}

\begin{proof}
By \eqref{eq:generator} and \eqref{eq:ampcert},
$\mathcal G_vV_q(z)=V_q(z)\sum_iz_i\Phi_v(q)_i/q_i\le0$ for every admitted $v$
and every $z$, including $z=0$. Localizing as in Theorem~\ref{thm:main} and
using $0\le V_q\le1$, the process $V_q(Z_t)$ is a bounded supermartingale under
any predictable policy, so $\E V_q(Z_t)\le V_q(z)=\prod_iq_i^{z_i}$. Extinction
by time $t$ implies $V_q(Z_t)=1$, whence
$\Prob_z(|Z_t|=0)\le\E V_q(Z_t)\le\prod_iq_i^{z_i}$. Letting $t\to\infty$ and
using that $0$ is absorbing gives the statement for eventual extinction. For a
continuation with $\Phi_{\mathrm{off}}(q)\le0$, apply the same supermartingale
bound to the continuation started from $Z_T$ and take expectations. Affineness
of $v\mapsto\Phi_v(q)=\PhiB(q)+vRq$ gives the last sentence.
\end{proof}

\begin{remark}\label{rem:ampvsexp}
Theorem~\ref{thm:amp} is the degenerate case $c=0$ of \eqref{eq:cert}, but it is
not a corollary of Theorem~\ref{thm:main} as stated there: condition
\eqref{eq:cert} with $c=0$ requires $Rq\le0$, which forces
$\Phi_v(q)\le\PhiB(q)\le0$ for all $v\ge0$ and is therefore a strictly stronger
requirement than \eqref{eq:ampcert} on a bounded action range. The two
certificates are genuinely different objects: \eqref{eq:cert} constrains the
direction $R$ relative to $1-q$ and yields a bound decaying in $B$;
\eqref{eq:ampcert} constrains the field at the endpoint of the action range and
yields a bound that does not decay at all.
\end{remark}

\begin{remark}[Cell-specific actions]\label{rem:cellwise}
Theorem~\ref{thm:amp} holds verbatim when each cell carries its own predictable
action $v\in\mathcal V$, because \eqref{eq:generator} sums per-cell contributions
and each is separately nonpositive. Theorem~\ref{thm:main}, by contrast, is a
statement about one common budget and must not be silently reinterpreted as a
sum of cell-specific exposures.
\end{remark}

\subsection{Sharpness of the amplitude threshold}
\label{sec:ampsharp}

Define the critical amplitude of a source as
\[
 v_c=\sup\bigl\{v\ge0:\ \exists\,q\in(0,1)^m\ \text{with}\
      \PhiB(q)\le0\ \text{and}\ \Phi_{v}(q)\le0\bigr\}.
\]
For $v_{\max}<v_c$, Theorem~\ref{thm:amp} gives a positive survival floor for
every policy. In the other direction, suppose the first-moment matrix
\eqref{eq:meanmatrix} at a constant admissible action $v_*$ has spectral
abscissa $\spb(A_{v_*})<0$.
Then, by Proposition~\ref{prop:weight} below, the constant action $v_*$ drives
the population to extinction almost surely, and the exposure required to reach
risk $\delta$ from $n$ founders is finite and of order $\log(n/\delta)$. In an
irreducible finite-type source with nondegenerate reproduction, $q^{(v)}<\one$
if and only if $\spb(A_v)>0$ \citep{athreya1972,kimmel2015}, and the minimal
extinction vector $q^{(v)}$ of the constant-$v$ source is itself a candidate in
the supremum defining $v_c$ whenever it also satisfies $\PhiB(q^{(v)})\le0$.
Consequently
\[
 v_c\ \le\ \inf\{v:\spb(A_v)<0\},
\]
and the two coincide whenever the minimal extinction vector at the larger
amplitude is also a baseline supersolution. Section~\ref{sec:memory} computes
$\inf\{v:\spb(A_v)<0\}$ exactly in the molecular source and exhibits amplitude
certificates at values just below it --- for instance at $e=.09$ when
$e_c=.092750\ldots$ at two sites, and at $e=.35$ when $e_c=.373695\ldots$ at
six --- while the construction is infeasible above it. We do not claim the two
quantities are equal, only that no gap is visible at the resolution computed.

\subsection{A separate finite-time floor}

\begin{proposition}[Finite-time floor]\label{prop:time}
For a binary source in which type changes preserve cell count and all controlled
death rates are at most $d_{\max}>0$,
\begin{equation}\label{eq:time}
 \Prob(|Z_T|>0)\ge1-\bigl(1-e^{-d_{\max}T}\bigr)^n,\qquad
 T\ge-\frac1{d_{\max}}\log\bigl[1-(1-\delta)^{1/n}\bigr]
\end{equation}
is necessary for finite-time risk at most $\delta$ from $n$ founders.
\end{proposition}

\begin{proof}
Let $x(t)=1-e^{-d_{\max}(T-t)}$ and $f(t,z)=1-x(t)^{|z|}$, with $f(t,0)=0$.
For $t<T$ and $k=|z|>0$, chemical changes contribute zero because they preserve
$|z|$, and
\[
 (\partial_t+\mathcal G)f
 =x^{k-1}(1-x)\Bigl(kd_{\max}-\sum_iz_id_i+x\sum_iz_ib_i\Bigr)\ \ge\ 0 .
\]
Localization and boundedness give the submartingale inequality. At $T$,
$f(T,z)=\one_{\{|z|>0\}}$ read as a polynomial value. Taking expectations proves
the first bound and rearrangement the second.
\end{proof}

The time and exposure restrictions are simultaneous necessary conditions; their
probability bounds are not multiplied. The time floor does not apply to eventual
survival after an arbitrarily long continuation. It requires non-death
reproduction to leave at least one descendant, as binary division does, and any
zero-offspring replacement event must be counted as an effective death.
